\chapter{Milestone M1.7: Generalized Bloch Supercells, Complexity Scaling, and Symmetry Stability}
\label{ch:m1_7_supercells}

% EDITORIAL QUESTION: How does the self-consistent mean-field Hamiltonian generalize to multi-sublattice supercells, and how is Brillouin zone folding represented algebraically?
\section{Algebraic Formulation of Honeycomb Supercells and BZ Folding}
\label{sec:m1_7_algebraic_supercells}

To investigate whether the two-sublattice ansatz suppresses incommensurate, stripy, or zigzag magnetic orders, Milestone M1.7 formulates the mean-field theory on an arbitrary $n_1 \times n_2$ supercell with $n_c = 2 n_1 n_2$ sublattices.

Let $\mathbf{a}_1, \mathbf{a}_2$ denote the primitive translation vectors of the honeycomb Bravais lattice with $|\mathbf{a}_1| = |\mathbf{a}_2| = 1$ and nearest-neighbor distance $d_{\mathrm{nn}} = 1/\sqrt{3}$. The supercell translation vectors are defined by:
\begin{equation}
\label{eq:m1_7_supercell_vectors}
\mathbf{A}_1 = n_1 \mathbf{a}_1, \qquad \mathbf{A}_2 = n_2 \mathbf{a}_2.
\end{equation}
The supercell contains $n_1 n_2$ primitive unit cells labeled by $(m_1, m_2)$ with $0 \le m_1 < n_1$ and $0 \le m_2 < n_2$. The $n_c = 2 n_1 n_2$ sublattices are partitioned into $A$-sites ($s = 2(m_1 + n_1 m_2)$) and $B$-sites ($s = 2(m_1 + n_1 m_2) + 1$).

\begin{theorem}[Brillouin Zone Folding and Fiber Expansion]
\label{thm:m1_7_bz_folding}
Let $\mathrm{BZ}_{\mathrm{prim}}$ and $\mathrm{BZ}_{\mathrm{super}}$ denote the first Brillouin zones of the primitive and supercell lattices, respectively. The volume ratio satisfies $\operatorname{Vol}(\mathrm{BZ}_{\mathrm{super}}) = \frac{1}{n_1 n_2} \operatorname{Vol}(\mathrm{BZ}_{\mathrm{prim}})$.
At each reduced wavevector $\mathbf{q} \in \mathrm{BZ}_{\mathrm{super}}$, the $4n_c \times 4n_c$ supercell Bloch Hamiltonian $\mathcal{H}_{\mathrm{super}}(\mathbf{q})$ is unitarily equivalent to the direct sum of primitive Hamiltonians evaluated at folded momenta:
\begin{equation}
\label{eq:m1_7_direct_sum_folding}
\mathcal{H}_{\mathrm{super}}(\mathbf{q}) \cong \bigoplus_{j_1=0}^{n_1-1} \bigoplus_{j_2=0}^{n_2-1} \mathcal{H}_{\mathrm{prim}}\left( \mathbf{q} + j_1 \mathbf{b}_1^{\mathrm{super}} + j_2 \mathbf{b}_2^{\mathrm{super}} \right).
\end{equation}
Consequently, the single-particle energy spectrum is identical, and the variational ground-state energy density satisfies:
\begin{equation}
\label{eq:m1_7_exact_energy_parity}
E_{\mathrm{site}}^{\mathrm{super}} \equiv \frac{1}{n_c} E_{\mathrm{cell}}^{\mathrm{super}} = \frac{1}{2} E_{\mathrm{cell}}^{\mathrm{prim}} \equiv E_{\mathrm{site}}^{\mathrm{prim}}.
\end{equation}
\end{theorem}
\begin{proof}
Let $\mathcal{U}$ denote the discrete Fourier transform matrix over the supercell translation group $G = \mathbb{Z}_{n_1} \times \mathbb{Z}_{n_2}$. Since the physical Kitaev exchange interactions and external Zeeman coupling are strictly invariant under primitive translations, the supercell Hamiltonian commutes with all elements of $G$. The characters of $G$ correspond to the folded reciprocal vectors $\mathbf{k}_{\mathbf{j}} = j_1 \mathbf{b}_1^{\mathrm{super}} + j_2 \mathbf{b}_2^{\mathrm{super}}$. Block diagonalization by $\mathcal{U}$ decomposes $\mathcal{H}_{\mathrm{super}}(\mathbf{q})$ into $n_1 n_2$ uncoupled $8 \times 8$ blocks $\mathcal{H}_{\mathrm{prim}}(\mathbf{q} + \mathbf{k}_{\mathbf{j}})$. Summing negative eigenvalues over $\mathbf{q} \in \mathrm{BZ}_{\mathrm{super}}$ reconstructs the full primitive BZ integral, guaranteeing exact equality of intensive energy densities.
\end{proof}

% EDITORIAL QUESTION: How does the computational cost of self-consistent mean-field evaluations scale with supercell size?
\section{Computational Complexity and Scaling Verification}
\label{sec:m1_7_complexity_scaling}

The evaluation map $\Phi$ requires diagonalizing $\mathcal{H}(\mathbf{q})$ at each point of the discrete $k$-mesh. For a supercell with $n_c$ sites, the matrix dimension is $M = 4n_c$, and the variational parameter space has dimension:
\begin{equation}
\label{eq:m1_7_dim_x}
D = \dim \mathcal{X} = 16 N_{\mathrm{bonds}} + 16 n_c + 3 n_c = 67 n_c.
\end{equation}
Under equivalent thermodynamic sampling, the number of supercell $k$-points is reduced to $N_k^{\mathrm{super}} = N_k^{\mathrm{prim}} / (n_1 n_2) \propto 1 / n_c$. The total eigensolver FLOP count per evaluation scales as:
\begin{equation}
\label{eq:m1_7_flop_scaling}
\mathcal{W}_{\Phi}(n_c) = N_k^{\mathrm{super}} \times \frac{16}{3} (4n_c)^3 = \frac{1024}{3} (N_k^{\mathrm{super}} n_c) \cdot n_c^2 \propto n_c^2.
\end{equation}

\begin{figure}[htbp]
\centering
\begin{tikzpicture}
\begin{axis}[
    labbook standard,
    xmode=log,
    ymode=log,
    log basis x={2},
    log basis y={10},
    xlabel={Supercell Size $n_c$ (Sublattices)},
    ylabel={Cost per Map Evaluation},
    legend pos=north west,
    xmin=1.5, xmax=45,
    ymin=0.5, ymax=600
]
    \addplot[labbookteal, thick, mark=*, mark size=2.5pt]
        table[x index=0, y index=7] {../../data/processed/benchmarks/m1_7/bloch_supercell_scaling.dat};
    \addlegendentry{FLOPs per Eval [MFLOP]}

    \addplot[labbookamber, thick, mark=square*, mark size=2pt]
        table[x index=0, y index=6] {../../data/processed/benchmarks/m1_7/bloch_supercell_scaling.dat};
    \addlegendentry{Wall Time per Eval [ms]}

    % Theoretical quadratic scaling guide: O(n_c^2)
    \addplot[labbookslate, dashed, thick, domain=2:32]
        {1.573 * (x/2)^2};
    \addlegendentry{Theoretical $\mathcal{O}(n_c^2)$ Folding Law}
\end{axis}
\end{tikzpicture}
\caption{Computational scaling of the self-consistent mean-field map evaluation across generalized honeycomb supercells $n_c \in \{2, 4, 8, 16, 32\}$ under equivalent thermodynamic sampling ($N_k \cdot n_c = 1152$). The measured computational work tracks the theoretical $\mathcal{O}(n_c^2)$ folded complexity law over two orders of magnitude, with intensive ground-state energy matching to within $5.55 \times 10^{-16}$.}
\label{fig:m1_7_supercell_scaling}
\end{figure}


\begin{proposition}[Empirical Supercell Scaling]
\label{prop:m1_7_empirical_scaling}
Numerical benchmarking across $n_c \in \{2, 4, 8, 16, 32\}$ sublattices at $h/|K| = 0.20$ (\texttt{scripts/m1\_7/benchmark\_bloch\_supercells.py}) establishes:
\begin{enumerate}
\item \textbf{Exact Intensive Energy Parity:} The ground-state energy per physical site satisfies $|E_{\mathrm{site}}(n_c) - E_{\mathrm{site}}(2)| \le 5.55 \times 10^{-16}$ across all sizes, matching primitive results to machine precision.
\item \textbf{Cubic Eigensolver Scaling:} The FLOP count per evaluation quadruples when $n_c$ doubles ($1.57 \to 6.29 \to 25.17 \to 100.66 \to 402.65\text{ MFLOPs}$), validating the $\mathcal{O}(n_c^2)$ folded complexity law with $R^2 > 0.99999$.
\item \textbf{Algorithmic Invariance:} The Joint DIIS solver converged in exactly $K = 23$ iterations across all five supercell sizes, demonstrating zero convergence degradation under supercell enlargement.
\end{enumerate}
\end{proposition}

% EDITORIAL QUESTION: Do unconstrained random perturbations trigger spontaneous translation or rotational symmetry breaking on supercells?
\section{Absence of Spontaneous Symmetry Breaking}
\label{sec:m1_7_symmetry_stability}

To verify whether the isotropic chiral spin liquid state is globally stable or merely constrained by initialization, we tested $N_{\mathrm{trials}} = 10$ unconstrained, symmetry-breaking random seed perturbations per supercell size for $n_c \in \{2, 4, 8, 16\}$ at $h/|K| = 0.20$.

\begin{theorem}[Ground-State Isotropy and Stability Against Modulation]
\label{thm:m1_7_stability}
Across all supercell sizes $n_c \in \{2, 4, 8, 16\}$:
\begin{enumerate}
\item \textbf{Universal Recovery:} Every unconstrained random trial converged to the true ground state ($P_{\mathrm{conv}} = 1.000, P_{\mathrm{iso}} = 1.000$).
\item \textbf{Vanishing Nematicity:} The residual spin polarization anisotropy satisfies:
\begin{equation}
N_{\mathrm{nem}} \equiv \frac{1}{n_c} \sum_{s=0}^{n_c-1} \left( |\sigma_s^x - \sigma_s^y| + |\sigma_s^y - \sigma_s^z| \right) \le 7.32 \times 10^{-8},
\end{equation}
confirming exact threefold $C_3$ rotational invariance.
\item \textbf{Absence of Plaquette Modulations:} Plaquette fluxes remain completely uniform across the lattice ($|W_p - W_{p'}| < 10^{-14}$), ruling out valence bond solids, stripy states, or staggered flux phases.
\end{enumerate}
\end{theorem}
